\documentclass[11pt]{scrartcl}
\usepackage[a4paper,left=2.4cm,right=3.2cm,top=2.4cm,bottom=2.4cm]{geometry}
\usepackage{../common/rlstudy}
\title{\textbf{Week 26: Extended reading}\\\large Causal biology, environment shifts, and task-aware enhancement}
\author{STUDY-RL}
\date{Sources checked October 5, 2026}
\begin{document}
\maketitle

\section*{How to use this reading pack}
Read the teaching guide first. These five papers cover different questions, and
their guarantees are not interchangeable. For each paper, identify the mathematical
object being recovered, its data assumptions, and the meaning of its evaluation.
The short excerpts below are verified source text. Commentary, discussion prompts,
and numerical thought experiments are supplied by this course.

The public Synlico event announcement motivates the topic. It does not supply a
reproducible technical description of the proprietary enhancement model. We therefore
include one independent task-driven augmentation paper and one verified paper by
Yingzhen Yang, without presenting either as the event's undisclosed algorithm.

\begin{keyidea}
Keep four records while reading:
\begin{enumerate}
\item the estimand or graph information the method targets;
\item which causal and observation mechanisms may change;
\item which assumptions are scientific inputs rather than learned outputs;
\item which components are absent from our small NumPy project.
\end{enumerate}
\end{keyidea}

For effect-estimation foundations, begin with Miguel A. Hern\'an and James M.
Robins (2020), \emph{Causal Inference: What If}, Chapman \& Hall/CRC.
\url{https://miguelhernan.org/whatifbook}. Chapters 1--3, 6--9, and 13 connect the
target question to consistency, exchangeability, positivity, graphs, measurement,
and standardization. The author page hosts the living online book and companion code.

This pack is an introduction to a research area, not sufficient preparation for
designing a biomedical causal study unaided. Its exercises are deliberately inspectable:
our covariances and effects are known because we choose the synthetic assignments.

\clearpage
\section*{Paper 1: multi-environment discovery with latent causes}
\textbf{Full citation.} Adam Li, Amin Jaber, Elias Bareinboim (2023).
\emph{Causal discovery from observational and interventional data across multiple
environments}. Advances in Neural Information Processing Systems 36, 16942--16956.
DOI: \texttt{10.52202/075280-0741}. No arXiv ID was verified; cite the official
proceedings version.
\url{https://proceedings.neurips.cc/paper_files/paper/2023/hash/368cba57d00902c752eaa9e4770bbbbe-Abstract-Conference.html}.

\paragraph{Summary.} The paper develops cross-domain structure-learning constraints
for systems that can contain hidden causes. Its selection-diagram and S-Markov
framework supports the S-FCI algorithm. The formal observational/interventional
equivalence is conditional on the modeled setting, rather than a claim that arbitrary
biological batches are experiments.

\paragraph{Why it matters for this week.} This is a public peer-reviewed connection
to Amin Jaber's multi-environment research. It explains why environments can supply
information beyond one observational distribution, while preserving unresolved structure.

\paragraph{Annotated key passage.}
\begin{quote}
``there exist latent confounders''
\end{quote}
\emph{Source location: abstract, proceedings PDF p.~1.}
\emph{Commentary.} This small phrase makes hidden common causes part of the problem,
rather than eliminating them by declaration. When reading any application, ask which
variables and context assumptions support its retained graph information.

\paragraph{Reading questions.} Write down the observed-variable set and the possible
latent causes. Locate the role of selection variables: which mechanisms can differ
between domains? Separate an observational domain from a domain with known laboratory
interventions. Explain what an unresolved orientation means for a proposed target effect.

\paragraph{Course-created thought experiment.} Our forward and reverse Gaussian SCMs
have covariance $\begin{pmatrix}1&0.8\\0.8&1\end{pmatrix}$. In one environment,
changing only the distribution of a measurement artifact need not distinguish their
biological directions. A controlled intervention on $X$ instead changes the response
mean to $0.8x$ in the forward model and leaves it zero in the reverse model. Explain
which additional intervention-specific assumptions make that contrast interpretable.

\paragraph{Implementation boundary.} Our project does not implement S-FCI,
conditional-independence tests, selection-diagram construction, or do-calculus.
It neither reproduces the paper's experiments nor inherits its soundness result.
Its standardization fixture assumes a measured confounder, a substantially narrower
effect-estimation setting than the paper's latent-variable discovery problem.

\begin{pitfall}
Replacing unknown endpoints by a favorite DAG and then simulating an intervention
does not remove the uncertainty in the endpoints. Track the alternatives that could
change the intervention answer.
\end{pitfall}

\clearpage
\section*{Paper 2: invariant causal prediction}
\textbf{Full citation.} Jonas Peters, Peter B\"uhlmann, Nicolai Meinshausen (2016).
\emph{Causal inference by using invariant prediction: identification and confidence
intervals}. Journal of the Royal Statistical Society, Series B 78(5), 947--1012.
DOI: \texttt{10.1111/rssb.12167}. arXiv: \texttt{1501.01332}.
\url{https://arxiv.org/abs/1501.01332}.

\paragraph{Summary.} ICP uses a target's invariant conditional mechanism across
environments to support parent-set inference with suitable testing assumptions.
In informative settings, invariance can help identify causal predictors; without
enough distinguishing variation, several sets can remain compatible.

\paragraph{Why it matters for this week.} It separates the target-parent problem
from recovery of an entire biological graph. Its formal testing machinery explains
why the project's residual summary must not be reported as causal confidence.

\paragraph{Annotated key passage.}
\begin{quote}
``the set of causal predictors becomes identifiable''
\end{quote}
\emph{Source location: Summary, journal p.~947; also the arXiv abstract.}
\emph{Commentary.} This statement concerns sufficient assumptions discussed by the
paper. In our work, ask which environment changes could distinguish predictor sets,
and which changes would instead violate the target's invariant mechanism.

\paragraph{Reconstruct the premise.} In the course equation
$Y=\beta^\top X_S+U_Y$, assume $\beta$ and the distribution of $U_Y$ remain fixed
and $U_Y$ stays independent of the predictors. Conditioning on $X_S=x$ gives the
same random quantity $\beta^\top x+U_Y$ in each environment. Changed predictor
marginals alone do not alter it. Now change $\beta$ or a measurement rule and identify
which equality no longer holds.

\paragraph{Diagnostic versus test.} The workbook compares residual means and variances.
Groups $(-1,1)$ and $(-2,2)$ have equal means and variances one and four. A mean-only
diagnostic misses their difference. Conversely, matching both moments would still
not establish the complete conditional-distribution premise.

\paragraph{Reading questions.} What is intersected, and why? What prevents a failure
to reject from becoming a proof of an individual edge? Which assumptions concern
interventions on the target? How are type-I error, power, and identification different?
Use the paper's formal statement when answering; our regression experiment supplies
none of those guarantees.

\begin{pitfall}
Ordinary held-out prediction error, an environment-average loss, and a calibrated
invariance test are different objects. A stable error bar does not turn one into another.
\end{pitfall}

\clearpage
\section*{Paper 3: a public example of task-driven augmentation}
\textbf{Full citation.} Krishna Chaitanya, Neerav Karani, Christian Baumgartner,
Olivio Donati, Anton Becker, Ender Konukoglu (2019; author order follows the cited
preprint). \emph{Semi-Supervised and Task-Driven Data Augmentation}.
Information Processing in Medical Imaging (IPMI 2019), LNCS 11492, 29--41.
DOI: \texttt{10.1007/978-3-030-20351-1\_3}. arXiv: \texttt{1902.05396}.
\url{https://arxiv.org/abs/1902.05396}.

\paragraph{Summary.} This medical-image segmentation work learns transformations
with a downstream task objective while maintaining image/annotation correspondence.
It supplies a concrete public meaning of task-driven enhancement. Its imaging results
do not establish biological graph recovery or an omics enhancement method.

\paragraph{Why it matters for this week.} Improving a downstream objective is a
different requirement from making transformed data look realistic. Label preservation
and a held-out evaluation protocol must be specified as part of the design.

\paragraph{Annotated key passage.}
\begin{quote}
``optimized to be conducive to the task''
\end{quote}
\emph{Source location: abstract, cited arXiv PDF p.~1.}
\emph{Commentary.} The task must be named. Our fixture's task is prediction MSE under
changing proxies, so a transformation selected for that task should not be reused
unexamined for estimating dependencies between genes.

\paragraph{Translate the idea carefully.} For segmentation, spatial changes must
keep the image and its mask consistent. For our continuous-marker fixture, permuting
the known nuisance changes a feature while preserving the outcome because the SCM
does not let that nuisance generate the outcome. If we permuted a true cause, the
same label-preservation argument would fail. Biological knowledge or a validated
observation model would be needed to justify transformations outside this fixture.

\paragraph{A design checklist for the exercise.} State which variable may change,
which target stays fixed, which distributions the transformation should preserve,
and which evaluation quantity will choose its strength. Fix the candidate family
before evaluation. Fit learned transformations only on permitted training patients.
Include the unchanged-data predictor and an ablation that omits the suspect feature.

\paragraph{Implementation boundary.} Our nuisance permutation is an instructor-created
heuristic. We do not reproduce this paper's generative networks, spatial transformations,
joint losses, image datasets, or experiments. We also do not attribute this independent
paper or our permutation to Yingzhen Yang or Synlico.

\begin{pitfall}
Augmented rows share their source information. Treating them as new independent
patients would misstate both sample size and uncertainty.
\end{pitfall}

\clearpage
\section*{Paper 4: negative controls for imputation artifacts}
\textbf{Full citation.} Tallulah S. Andrews and Martin Hemberg (2018; version 2
published 2019). \emph{False signals induced by single-cell imputation}.
F1000Research 7:1740. DOI: \texttt{10.12688/f1000research.16613.2}.
No arXiv ID was verified.
\url{https://f1000research.com/articles/7-1740}.

\paragraph{Summary.} The study evaluates false signals using negative controls
after single-cell imputation. It motivates examining manufactured dependencies,
not only recovery of known positive signals. This is empirical evidence about
tested settings, not a universal theorem that every imputation harms every task.

\paragraph{Why it matters for this week.} Discovery depends on statistical relationships.
An enhancement pipeline can change those relationships, so reconstruction quality
or a visually smooth embedding is not a sufficient structure-learning evaluation.

\paragraph{Annotated key passage.}
\begin{quote}
``all methods can introduce false positive signals into data''
\end{quote}
\emph{Source location: Introduction, paragraph discussing negative controls in the
version-2 primary article.}
\emph{Commentary.} The wording expresses a possibility in the study's evaluated
methods and settings. In our lesson, use an explicitly absent dependency as a control
and state exactly which transformation and sample it diagnoses.

\paragraph{Course-created counterexample.} Let $G_1,G_2$ be independent, mean-zero
markers with variance one. Define $H_1=H_2=(G_1+G_2)/2$. Their variances and covariance
are each $1/2$, giving correlation one. This transparent construction is not a
reimplementation of a published method. It demonstrates why shared smoothing can
alter dependencies even when output variance decreases.

\paragraph{Missingness versus zero.} The project's NaN mask is known and MCAR.
Zero is retained as an observed value. In real measurements, a zero does not by itself
reveal a missing-value mechanism. Before borrowing an imputation result, record the
assay, normalization, observation mask, and unit of independence. Otherwise the toy
missingness contract can silently become an unjustified biological assumption.

\paragraph{Reading questions.} What constitutes a negative control in the study?
Which downstream inference is being tested? What changes if donor-related observations
appear in both fitting and evaluation? How could a pipeline improve prediction and
still create false candidate edges? Identify separate metrics for those two outcomes.

\begin{pitfall}
A marker value created by a model is not an additional experiment on that marker.
Track raw and enhanced observations separately when assessing evidence.
\end{pitfall}

\clearpage
\section*{Paper 5: verified research context for Yingzhen Yang}
\textbf{Full citation.} Yingzhen Yang and Ping Li (2022).
\emph{Noisy $\ell_0$-sparse subspace clustering on dimensionality reduced data}.
Proceedings of the Thirty-Eighth Conference on Uncertainty in Artificial Intelligence
(UAI 2022), PMLR 180:2235--2245. arXiv: \texttt{2206.11079}.
\url{https://proceedings.mlr.press/v180/yang22e.html}.

\paragraph{Summary.} The paper studies sparse subspace detection with noisy data
and dimensionality reduction under explicit mathematical models. It provides related
research context about learning structure despite noise. Subspace detection is not
the same object as identifying a directed causal graph.

\paragraph{Why it matters for this week.} The speaker's verified public work makes
noise and structural assumptions a useful discussion topic. It does not disclose the
particular enhancement model referenced by the Synlico event announcement.

\paragraph{Annotated key passage.}
\begin{quote}
``Real data often suffer from noise''
\end{quote}
\emph{Source location: abstract, proceedings PDF p.~1.}
\emph{Commentary.} A theorem's treatment of noise specifies more than the word
``noisy.'' Ask about the disturbance distribution, geometric conditions, observation
mechanism, and scale at which a guarantee applies before transferring it to patient data.

\paragraph{Separate structural objects.} A clustering relation can summarize proximity
or membership in a low-dimensional subspace. A causal graph asserts directed mechanism
relationships. An accurate cluster can be useful in a biological workflow without
identifying the direction of a gene effect. Conversely, a causal argument does not
automatically prove that a dimensionality-reduction step preserves every needed signal.

\paragraph{Course-created translation exercise.} In the classical measurement example
$W=X+\epsilon$, $Y=2X+U$, assume independent mean-zero disturbances with
$\Var(X)=\Var(\epsilon)=1$. The population regression coefficient on $W$ is
$2\Var(X)/(\Var(X)+\Var(\epsilon))=1$. Reducing measurement noise could change
this coefficient without identifying whether $X$ causes $Y$ in observationally
ambiguous settings. Keep the measurement and graph-identification questions separate.

\paragraph{Questions for the speaker.} What is the downstream task for enhancement?
Which mathematical noise model describes the observed assay? What changes when the
training and evaluation units are donors rather than cells? Which controls reveal
invented structure? Which guarantees apply to the enhancement implementation actually
used? These questions are invitations to clarify, not assertions about undisclosed methods.

\begin{pitfall}
A verified publication by a speaker does not establish that it is the method used
in a company platform or the method presented at a particular event.
\end{pitfall}

\clearpage
\section*{Connecting the papers to an evidence-backed experiment}
Our project deliberately separates an identifiable measured-confounder ATE fixture
from a prediction experiment with changing proxies. It includes patient grouping,
frozen preprocessing, validation-only candidate selection, a no-proxy ablation, and
a fixed-model patient bootstrap. It does not merge these into a claim about real
drug efficacy or recovery of biological gene networks.

The central numerical findings are easy to inspect: adjusting the known binary
confounder recovers an effect near one; the unadjusted difference is near 2.2; a
validation-selected transformation improves held-out prediction in the default fixture;
removing the proxy entirely is often better in the unseen environment. The final
point explains why enhancement success should be reported with its baselines and limits.

\paragraph{Broader multi-context background.} Joris M. Mooij, Sara Magliacane,
Tom Claassen (2020), \emph{Joint Causal Inference from Multiple Contexts}, JMLR
21(99), 1--108, arXiv:\texttt{1611.10351}.
\url{https://jmlr.org/papers/v21/17-123.html}.
Use this framework reference to distinguish context assumptions from particular
PC/FCI implementations. It is supplementary reading rather than a sixth annotated paper.

\paragraph{A route to experimental validation.} Replogle and colleagues (2022),
\emph{Mapping information-rich genotype-phenotype landscapes with genome-scale
Perturb-seq}, Cell 185(14), 2559--2575.e28, DOI:\texttt{10.1016/j.cell.2022.05.013}.
\url{https://pmc.ncbi.nlm.nih.gov/articles/PMC9380471/}.
This primary perturbation study illustrates molecular hypothesis testing. Its full
author list is recorded in the source dossier. A cell-context perturbation response
is distinct from a patient's treatment effect.

\paragraph{Your one-page memo.} Name the question, unit, population, intervention,
and timing. List the assumptions needed for identification and across-environment
transport. Show the exact splits and fit/select/evaluate boundaries. Report baseline,
enhancement, and ablation on one held-out metric, then identify one unresolved causal
alternative and one experiment that could discriminate it. Distinguish a conditional
sampling interval from uncertainty about the causal model.

\paragraph{Sourced from.}
\href{../research/causal_biology.md}{\texttt{research/causal\_biology.md}}
records source URLs, versions, event provenance, overlap with existing weeks, and
the distinction between researched methods and the course-created mini-project.
The related support and matching background is in
\href{../research/bandits_matching_sota.md}{\texttt{research/bandits\_matching\_sota.md}}.
\end{document}
